% =====================================================================
%  main.tex - Slide bảo vệ bài tập lớn (khung: theme MLWNet của người làm đồ án)
%  Font: config/fonts.tex | Logo, ảnh, footer: config/settings.tex
%  Giao diện: theme/beamerthemeMLWNet.sty
%  Biên dịch: XeLaTeX -> BibTeX -> XeLaTeX -> XeLaTeX (Overleaf: Compiler = XeLaTeX)
%  Số liệu lấy từ báo cáo (report/overleaf), không gõ số mới.
% =====================================================================
\makeatletter\def\input@path{{config/}{theme/}}\makeatother
% =====================================================================
%  config/fonts.tex - TINH CHỈNH FONT
%  File này chỉ khai báo giá trị; theme/beamerthemeMLWNet.sty đọc và
%  áp dụng. Được \input TRƯỚC \documentclass nên chỉ dùng \newcommand.
% =====================================================================

% ---------------------------------------------------------------------
% 1. Họ font chữ (cần XeLaTeX hoặc LuaLaTeX)
% ---------------------------------------------------------------------
%  - mlwFontPath  : thư mục chứa .ttf/.otf (có "/" ở cuối).
%                   Để TRỐNG {} nếu muốn dùng font đã cài trong hệ thống.
%  - mlwFontName  : tiền tố tên file, VD BeVietnamPro -> BeVietnamPro-Regular.ttf
%                   (hoặc tên font hệ thống nếu mlwFontPath trống).
%  - mlwFontExt   : phần mở rộng (.ttf hoặc .otf).
%  - Hậu tố từng kiểu chữ; file nào thiếu thì theme tự bỏ qua kiểu đó.
\newcommand{\mlwFontPath}{fonts/}
\newcommand{\mlwFontName}{BeVietnamPro}
\newcommand{\mlwFontExt}{.ttf}
\newcommand{\mlwFontRegular}{-Regular}
\newcommand{\mlwFontBold}{-Bold}
\newcommand{\mlwFontItalic}{-Italic}
\newcommand{\mlwFontBoldItalic}{-BoldItalic}

%  Tỉ lệ co giãn (Be Vietnam Pro khá rộng; 0.95 giúp chữ gọn hơn)
\newcommand{\mlwFontScale}{0.95}

%  Ví dụ đổi sang font hệ thống (bỏ % để dùng):
% \renewcommand{\mlwFontPath}{}
% \renewcommand{\mlwFontName}{Arial}

% ---------------------------------------------------------------------
% 2. Font toán
% ---------------------------------------------------------------------
%  - cm    : Computer Modern có chân (mặc định LaTeX, an toàn nhất)
%  - sans  : chữ toán không chân theo font chính (kiểu Beamer gốc)
%  - Tên một font OpenType Math, dùng unicode-math, VD:
%      Latin Modern Math | TeX Gyre Pagella Math | TeX Gyre Termes Math
%  Ghi tên file (latinmodern-math.otf) thay cho tên font: Tectonic không tìm font theo tên hệ thống
\newcommand{\mlwMathFont}{latinmodern-math.otf}

% ---------------------------------------------------------------------
% 3. Cỡ chữ
% ---------------------------------------------------------------------
%  Cỡ chữ gốc: 8pt, 9pt, 10pt, 11pt, 12pt, 14pt, 17pt, 20pt
\newcommand{\mlwBaseSize}{10pt}

%  Cỡ chữ từng thành phần (\tiny \scriptsize \footnotesize \small
%  \normalsize \large \Large \LARGE \huge \Huge)
\newcommand{\mlwSizeFrameTitle}{\large}        % tiêu đề slide
\newcommand{\mlwSizeTitle}{\Large}             % tiêu đề trang bìa
\newcommand{\mlwSizeSection}{\LARGE}           % slide chuyển mục
\newcommand{\mlwSizeFooter}{\scriptsize}       % dòng chân trang
\newcommand{\mlwSizeSource}{\scriptsize}       % chú thích nguồn \src{}
\newcommand{\mlwSizeBib}{\scriptsize}          % danh mục tài liệu tham khảo
\newcommand{\mlwSizeQA}{\Huge}                 % chữ "HỎI ĐÁP"

% =====================================================================
%  config/settings.tex - THIẾT LẬP CHUNG
%  Logo, thư mục ảnh, dòng footer và các công tắc hiển thị.
%  Được \input TRƯỚC \documentclass nên chỉ dùng \newcommand.
% =====================================================================

% ---------------------------------------------------------------------
% 1. Khổ slide
% ---------------------------------------------------------------------
%  169 (16:9) | 1610 (16:10) | 43 (4:3) | 32 (3:2)
\newcommand{\mlwAspectRatio}{169}

% ---------------------------------------------------------------------
% 2. Hình ảnh
% ---------------------------------------------------------------------
%  Thư mục ảnh (có "/" ở cuối). \includegraphics và \figorbox
%  đều tìm ảnh tại đây, nên trong main.tex chỉ cần ghi tên file.
\newcommand{\mlwImageDir}{images/}

%  Logo trường (nằm trong \mlwImageDir). Thiếu file -> tự ẩn logo.
\newcommand{\mlwLogoFile}{logo.jpg}
\newcommand{\mlwLogoHeaderHeight}{1.2cm}       % logo góc phải tiêu đề slide
\newcommand{\mlwLogoTitleHeight}{1.7cm}        % logo trên trang bìa

%  true: logo ở mọi slide | false: chỉ ở trang bìa
\newcommand{\mlwLogoOnEveryFrame}{true}

% ---------------------------------------------------------------------
% 3. Dòng footer
% ---------------------------------------------------------------------
%  Chữ nghiêng bên trái chân trang. Để trống {} sẽ dùng \title ngắn
%  (\title[ngắn]{dài}) hoặc tiêu đề đầy đủ.
\newcommand{\mlwFooterText}{Phát hiện, phân loại và theo dõi phương tiện giao thông -- Dương Minh Chiến}

%  true: "5 / 20" | false: chỉ "5"
\newcommand{\mlwShowTotalFrames}{true}

% ---------------------------------------------------------------------
% 4. Bố cục và nhãn
% ---------------------------------------------------------------------
%  Tự chèn slide chuyển mục ở mỗi \section
\newcommand{\mlwSectionDivider}{true}

%  Nhãn trên trang bìa
\newcommand{\mlwLabelPresenter}{Người trình bày}
\newcommand{\mlwLabelStudentID}{MSSV}
\newcommand{\mlwLabelAdvisor}{Giảng viên hướng dẫn}
\newcommand{\mlwLabelQA}{HỎI ĐÁP}
\newcommand{\mlwLabelMissingFig}{Chưa có tệp hình:}

% ---------------------------------------------------------------------
% 5. Tài liệu tham khảo (BibTeX)
% ---------------------------------------------------------------------
%  File .bib (không ghi đuôi) và kiểu trích dẫn. Nếu máy không có
%  IEEEtran.bst, theme tự chuyển sang ieeetr (có sẵn trong TeX Live).
\newcommand{\mlwBibFile}{references}
\newcommand{\mlwBibStyle}{IEEEtran}

\documentclass[aspectratio=\mlwAspectRatio, \mlwBaseSize, t]{beamer}
\usetheme{MLWNet}
\usepackage{tabularx}
\newcolumntype{Y}{>{\raggedright\arraybackslash}X}
% Cột chữ căn trái có độ rộng cố định: cột p{} mặc định dàn đều, ô hẹp bị khoảng trắng lớn giữa các từ
\newcolumntype{P}[1]{>{\raggedright\arraybackslash}p{#1}}

% ---------------------------------------------------------------------
% Thông tin trang bìa
% ---------------------------------------------------------------------
\title[Phát hiện, phân loại và theo dõi phương tiện giao thông]{Ứng dụng học sâu trong phát hiện, phân loại\\
  và theo dõi phương tiện giao thông}
\sourceline{Bài tập lớn môn Học sâu}
\presenter{Dương Minh Chiến}
\studentid{12423062}
\advisor{PGS.TS. Nguyễn Minh Tiến}
\institute{Khoa Công nghệ Thông tin - Trường Đại học Công nghệ Kỹ thuật Hưng Yên}
\date{Hưng Yên, 2026}

% Ô số liệu kèm khoảng tin cậy: giá trị ở trên, [a–b] chữ nhỏ màu xám ở dưới
\newcommand{\ci}[2]{{\renewcommand{\arraystretch}{0.9}%
  \begin{tabular}[t]{@{}c@{}}#1\\{\tiny\color{MidGray}[#2]}\end{tabular}}}
% Tiêu đề cột hai dòng (chữ đậm), để cột không phải rộng theo tiêu đề dài
\newcommand{\hd}[2]{{\renewcommand{\arraystretch}{0.95}%
  \begin{tabular}[b]{@{}c@{}}\textbf{#1}\\\textbf{#2}\end{tabular}}}

\begin{document}

\utehytitlepage

\begin{frame}{Tiến Trình Trình Bày}
  \vspace{0.6em}
  \tableofcontents
\end{frame}

% =====================================================================
\section{Giới thiệu bài toán}
% =====================================================================

\subsection{Bối cảnh}
\begin{frame}{\stitle}
  \begin{columns}[T]
    \begin{column}{0.50\textwidth}
      \begin{itemize}
        \item Việt Nam có \alert{77 triệu xe máy} đăng ký (9/2024)\ct{tienphong2024xemay}.
        \item Xe máy đi chung với ô tô, buýt, tải; ít giữ làn, dồn thành cụm khi chờ đèn đỏ.
        \item Hà Nội vận hành trung tâm điều khiển giao thông thông minh (12/2025) với \alert{1.837 camera} tại 195 nút
              giao\ct{vneconomy2025camera}.
        \item Không thể có người xem từng camera: máy phải tự \alert{tìm xe, biết loại xe, nhận ra cùng một xe} qua
              các frame.
      \end{itemize}
    \end{column}
    \begin{column}{0.46\textwidth}
      \centering
      \figorbox{intro_day.jpg}{width=\linewidth,height=0.30\textheight,keepaspectratio}{2.3cm}
      \vspace{0.15em}
      \figorbox{intro_night.jpg}{width=\linewidth,height=0.30\textheight,keepaspectratio}{2.3cm}
      \src{Nguồn video: Raw Street Capture 101, Midnight Soundscape.}
    \end{column}
  \end{columns}
\end{frame}

\subsection{Bài toán}
\begin{frame}{\stitle}
  \begin{columns}[T]
    \begin{column}{0.44\textwidth}
      \begin{block}{Ba bài toán}
        \textbf{Phát hiện}: box của từng xe.\\
        \textbf{Phân loại}: xe máy, ô tô, buýt, tải.\\
        \textbf{Theo dõi}: một ID cho mỗi xe qua các frame.
      \end{block}
      \begin{block}{Đầu vào và đầu ra}
        \textbf{Vào:} video camera cố định nhìn chéo từ trên cao, ngày và đêm.\\
        \textbf{Ra:} box, lớp, ID; ứng dụng đếm xe qua vạch.
      \end{block}
    \end{column}
    \begin{column}{0.54\textwidth}
      \centering
      \figorbox{demo_pvd_goc_a.jpg}{width=\linewidth,height=0.58\textheight,keepaspectratio}{4cm}
      \src{Kết quả của hệ thống trên clip Phạm Văn Đồng (nguồn video: Vietnam Naturally).}
    \end{column}
  \end{columns}
\end{frame}

\subsection{Khó khăn}
\begin{frame}{\stitle}
  \begin{columns}[T]
    \begin{column}{0.48\textwidth}
      \begin{itemize}
        \item \textbf{Mật độ xe máy rất cao}: 2.982 xe máy trên 30 frame test Hà Nội, che khuất lẫn nhau.
        \item \textbf{Xe rất nhỏ ở xa}: 52\% xe ở cảnh đại lộ cao dưới 32 điểm ảnh.
        \item \textbf{Ban đêm}: đèn pha chói, xe tối lẫn nền, xe xa chỉ thấy đèn.
      \end{itemize}
    \end{column}
    \begin{column}{0.48\textwidth}
      \begin{itemize}
        \item \textbf{Lệch miền}: mô hình COCO gốc chỉ đạt mAP50 0.686 trên test Hà Nội.
        \item \textbf{Thiếu dữ liệu đúng góc camera giao thông}; định nghĩa lớp xe khác nhau giữa các bộ dữ liệu.
        \item \textbf{Thời gian thực}: video 30 frame/s, nhiều camera.
      \end{itemize}
    \end{column}
  \end{columns}
\end{frame}

\subsection{Mục tiêu và phạm vi}
\begin{frame}{\stitle}
  \begin{block}{Mục tiêu}
    Hệ thống phát hiện, phân loại và theo dõi phương tiện từ video camera giao thông Việt Nam có
    \alert{khả năng triển khai thực tế}: đủ chính xác trên camera chưa dùng khi huấn luyện, cả ngày lẫn đêm, và chạy
    kịp thời gian thực trên máy chủ có GPU.
  \end{block}
  \begin{block}{Phạm vi}
    \begin{itemize}
      \item Camera cố định nhìn chéo từ trên cao, kiểu camera giao thông; ban ngày và ban đêm.
      \item Bốn lớp: \texttt{motorbike}, \texttt{car}, \texttt{bus}, \texttt{truck}; box xe máy bao cả người lái.
      \item Đánh giá phát hiện bằng mAP (COCO), theo dõi bằng HOTA, MOTA, IDF1 (MOTChallenge).
    \end{itemize}
  \end{block}
\end{frame}

\subsection{Nội dung và nguyên tắc đánh giá}
\begin{frame}{\stitle}
  \begin{columns}[T]
    \begin{column}{0.48\textwidth}
      \begin{block}{Nội dung thực hiện}
        \begin{itemize}
          \item Gộp dữ liệu, tự xây tập test Việt Nam có người kiểm.
          \item Tinh chỉnh YOLO26 qua nhiều phiên bản dữ liệu.
          \item So sánh CNN (YOLO26) với Transformer (RT-DETRv2-S).
          \item Chọn tracker, đếm xe qua vạch, web demo.
        \end{itemize}
      \end{block}
    \end{column}
    \begin{column}{0.48\textwidth}
      \begin{block}{Nguyên tắc đánh giá}
        \begin{itemize}
          \item Mọi lựa chọn chốt trên tập \alert{dev}, theo tiêu chí viết trước khi chạy.
          \item Tập \alert{test} chỉ để báo cáo kết quả.
          \item Chênh dưới 0.005 coi là hoà; hoà thì chọn bản nhanh hơn.
        \end{itemize}
      \end{block}
    \end{column}
  \end{columns}
\end{frame}

% =====================================================================
\section{Cơ sở lý thuyết}
% =====================================================================

\subsection{YOLO26: phát hiện một giai đoạn, không cần NMS}
\begin{frame}{\stitle}
  \begin{columns}[T]
    \begin{column}{0.56\textwidth}
      \centering
      \figorbox{paper_yolo26_arch.png}{width=\linewidth,height=0.64\textheight,keepaspectratio}{4.6cm}
      \src{Nguồn: \cite{jocher2026yolo26}.}
    \end{column}
    \begin{column}{0.42\textwidth}
      \begin{itemize}
        \item Họ YOLO\ct{redmon2016yolo}: một lần đưa ảnh qua mạng cho ra mọi box và lớp.
        \item \textbf{Backbone} trích đặc trưng, \textbf{neck} trộn nhiều mức, \textbf{head} dự đoán ở 1/8, 1/16, 1/32.
        \item YOLO26: đầu một-một khi suy luận nên \alert{không cần NMS}.
        \item Hai cỡ: \textbf{YOLO26n} (2.5 M tham số), \textbf{YOLO26s} (10.0 M).
      \end{itemize}
    \end{column}
  \end{columns}
\end{frame}

\subsection{RT-DETRv2: phát hiện bằng Transformer}
\begin{frame}{\stitle}
  \begin{columns}[T]
    \begin{column}{0.56\textwidth}
      \centering
      \figorbox{paper_rtdetr.png}{width=\linewidth,height=0.60\textheight,keepaspectratio}{4.6cm}
      \src{Nguồn: \cite{zhao2024rtdetr}.}
    \end{column}
    \begin{column}{0.42\textwidth}
      \begin{itemize}
        \item DETR\ct{carion2020detr}: phát hiện là dự đoán tập hợp, ghép một-một bằng Hungarian.
        \item RT-DETR: bộ mã hoá lai (AIFI + CCFF), chọn truy vấn theo độ bất định, attention biến dạng.
        \item RT-DETRv2-S\ct{lv2024rtdetrv2}: ResNet-18, 20.1 M tham số.
        \item Ưu: box khít. Nhược: \alert{nặng hơn YOLO}.
      \end{itemize}
    \end{column}
  \end{columns}
\end{frame}

\subsection{Theo dõi đa đối tượng}
\begin{frame}{\stitle}
  % Hình ByteTrack gần vuông đặt cạnh chữ; sơ đồ BoT-SORT rất dẹt (4.6 : 1) nên trải hết bề ngang ở dưới cho đọc được
  \begin{columns}[T]
    \begin{column}{0.56\textwidth}
      \begin{itemize}
        \item \textbf{SORT}\ct{bewley2016sort}: Kalman dự đoán vị trí, Hungarian ghép box theo IoU.
        \item \textbf{ByteTrack}\ct{zhang2022bytetrack}: ghép thêm lần hai với box điểm thấp, xe bị che không mất
              track.
        \item \textbf{BoT-SORT}\ct{aharon2022botsort}: Kalman theo rộng và cao box, bù chuyển động camera, tuỳ chọn
              ReID.
        \item Ngưỡng tin cậy detector đặt \alert{0.1} để có box điểm thấp cho lần ghép hai.
      \end{itemize}
    \end{column}
    \begin{column}{0.42\textwidth}
      \centering
      \figorbox{paper_bytetrack.jpg}{width=\linewidth,height=0.40\textheight,keepaspectratio}{3cm}
      \src{ByteTrack. Nguồn: \cite{zhang2022bytetrack}.}
    \end{column}
  \end{columns}
  \vspace{0.1em}
  \centering
  \figorbox{paper_botsort.png}{width=\textwidth,height=0.22\textheight,keepaspectratio}{2.4cm}
  \src{Luồng xử lý của BoT-SORT. Nguồn: \cite{aharon2022botsort}.}
\end{frame}

\subsection{Chỉ số đánh giá}
\begin{frame}{\stitle}
  \begin{columns}[T]
    \begin{column}{0.48\textwidth}
      \begin{block}{Phát hiện: AP, mAP (COCO\ct{lin2014coco})}
        \[
          \mathrm{AP} = \frac{1}{101} \sum_{r \in \{0, 0.01, \ldots, 1\}} p_{\mathrm{interp}}(r)
        \]
        \textbf{mAP50 chính}: trung bình AP50 của xe máy và ô tô (buýt, tải quá ít mẫu trong tập test).
      \end{block}
    \end{column}
    \begin{column}{0.48\textwidth}
      \begin{block}{Theo dõi: MOTA, IDF1, HOTA}
        \[
          \mathrm{MOTA} = 1 - \frac{\sum_t (\mathrm{FN}_t + \mathrm{FP}_t + \mathrm{IDSW}_t)}{\sum_t \mathrm{GT}_t}
        \]
        \[
          \mathrm{HOTA}_\alpha = \sqrt{\mathrm{DetA}_\alpha \cdot \mathrm{AssA}_\alpha}
        \]
        IDF1, IDR: giữ đúng danh tính\ct{ristani2016idf1}; HOTA tách phát hiện và liên kết\ct{luiten2021hota}.
      \end{block}
    \end{column}
  \end{columns}
\end{frame}

% =====================================================================
\section{Dữ liệu và xây dựng hệ thống}
% =====================================================================

\subsection{Dữ liệu huấn luyện: gộp bảy nguồn}
\begin{frame}{\stitle}
  \centering
  \small
  \begin{tabularx}{\linewidth}{@{}P{4.1cm}Yr@{}}
    \toprule
    \textbf{Nguồn} & \textbf{Đặc điểm} & \textbf{Ảnh train} \\
    \midrule
    vietnamese-vehicle\ct{roboflow_vnvehicle} & Xe Việt Nam nhiều cảnh; nguồn duy nhất chia cả val, test & 929 \\
    cantho, vehicle-detection & Gần như một camera mỗi bộ; chỉ dùng train & 300 + 300 \\
    VisDrone\ct{zhu2022visdrone} & Rất nhiều xe máy nhỏ nhìn từ trên cao & 909 \\
    UVH-26\ct{sharma2025uvh26} & Đúng góc camera giao thông, khoảng 2.800 camera & 2.000 \\
    Video ngã tư Hà Nội & Đúng cảnh mục tiêu, nhãn tự động, lặp 4 lần & 69 $\times$ 4 \\
    \rowcolor{HiRow}
    Video TP.HCM ban đêm & Hai camera đêm, nhãn tự động (từ v6) & 288 \\
    \bottomrule
  \end{tabularx}\par
  \vspace{0.6em}
  \begin{itemize}
    \item Bộ dữ liệu v6: \alert{5.002 ảnh train}, 320 ảnh val, 308 ảnh test.
    \item Không dùng cách chia sẵn của Roboflow: ảnh gần trùng được gom nhóm bằng dHash rồi chia theo nhóm, tránh test
          cùng cảnh với train.
  \end{itemize}
\end{frame}

\subsection{Các tập đánh giá: đo trên camera mới}
\begin{frame}{\stitle}
  \begin{columns}[T]
    \begin{column}{0.46\textwidth}
      \begin{block}{Dev (dùng để chọn)}
        \begin{itemize}
          \item Val Hà Nội: 10 frame, nhãn người.
          \item Dev đêm: 32 frame nút giao Gia Định.
          \item VnTraffic\ct{le2026vntraffic} và clip dev Phạm Văn Đồng góc C: chọn tracker.
        \end{itemize}
      \end{block}
    \end{column}
    \begin{column}{0.50\textwidth}
      \begin{block}{Test (chỉ để báo cáo)}
        \begin{itemize}
          \item Test Hà Nội 30 frame: \alert{cùng camera} với frame train.
          \item \texttt{vn\_extra} ngày, đêm: \alert{camera mới}, 30 + 30 frame.
          \item Clip Phạm Văn Đồng 20 giây, 183 quỹ đạo nhãn tay.
          \item 176 lượt xe đếm tay ở hai clip hầm chui đêm.
          \item Val UVH-26, UA-DETRAC\ct{wen2020uadetrac}.
        \end{itemize}
      \end{block}
    \end{column}
  \end{columns}
\end{frame}

\subsection{Gán nhãn dữ liệu}
\begin{frame}{\stitle}
  \begin{columns}[T]
    \begin{column}{0.48\textwidth}
      \begin{block}{Quy trình ba bước}
        Mô hình tạo nhãn sơ bộ $\rightarrow$ người gán sửa từng box $\rightarrow$ mô hình khác chỉ ra chỗ có thể sót.
      \end{block}
      \begin{block}{Quy tắc}
        \begin{itemize}
          \item Box xe máy bao cả người lái.
          \item Xe mơ hồ gán lớp \texttt{ignore} (như KITTI, COCO).
          \item Van nhỏ là ô tô, cỡ Ford Transit là buýt.
        \end{itemize}
      \end{block}
    \end{column}
    \begin{column}{0.48\textwidth}
      \begin{block}{Nhãn tự động cho frame train}
        Ghép nhiều mô hình ``thầy'' bằng WBF\ct{solovyev2021wbf} và SAHI\ct{akyon2022sahi}; val Hà Nội: P 0.905,
        R 0.853.
      \end{block}
      \begin{block}{Tránh thiên lệch}
        Nhãn sơ bộ của tập test mới lấy từ YOLO26x COCO, \alert{không phải mô hình đang so sánh}.
      \end{block}
    \end{column}
  \end{columns}
\end{frame}

\subsection{Quá trình cải tiến mô hình}
\begin{frame}{\stitle}
  \centering
  \small
  \begin{tabularx}{\linewidth}{@{}lYccc@{}}
    \toprule
    \textbf{Bản} & \textbf{Thay đổi} & \textbf{Ảnh} & \hd{Test Hà Nội}{mAP50} & \hd{\texttt{vn\_extra} đêm}{mAP50} \\
    \midrule
    COCO gốc & YOLO26n chưa tinh chỉnh & 1280 & 0.686 & 0.437 \\
    v1 & Roboflow + VisDrone & 640 & 0.395 & 0.183 \\
    v2 & Ảnh 960, VisDrone nhiều xe máy & 960 & 0.890 & 0.496 \\
    v3 & + frame Hà Nội nhãn tự động & 960 & 0.921 & 0.437 \\
    v4 & Như v3, ảnh 1280 & 1280 & 0.919 & 0.481 \\
    v5 & + 2.000 ảnh UVH-26 (góc camera giao thông) & 960 & \textbf{0.923} & 0.628 \\
    v6 & + 288 frame đêm (\alert{bản nhẹ}, YOLO26n) & 960 & 0.921 & 0.755 \\
    \rowcolor{HiRow}
    YOLO26s-v6 & YOLO26s học dữ liệu v6 (\alert{mô hình chính}) & 960 & 0.920 & \textbf{0.887} \\
    \bottomrule
  \end{tabularx}\par
  \vspace{0.5em}
  \begin{itemize}
    \item Dữ liệu \alert{đúng góc camera} (v5) và \alert{đúng cảnh đêm} (v6) có tác dụng lớn hơn tăng ảnh lên 1280 (v4).
  \end{itemize}
\end{frame}

\subsection{Tiêu chí chọn mô hình viết trước khi chạy}
\begin{frame}{\stitle}
  \centering
  \small
  \begin{tabularx}{\linewidth}{@{}Ycc@{}}
    \toprule
    \textbf{Điều kiện trên tập dev (cả hai lần huấn luyện phải đạt)} & \hd{v6}{(seed 42 / 1)} & \hd{YOLO26s-v6}{(seed 42 / 1)} \\
    \midrule
    mAP50 dev đêm cao hơn lần tốt nhất của bản cũ ($> 0.499$; $> 0.641$) & 0.599 / 0.618 & 0.711 / 0.712 \\
    mAP50 val Hà Nội không kém quá 0.005 ($\ge 0.909$) & 0.914 / 0.914 & 0.911 / 0.916 \\
    IDR VnTraffic không kém quá 0.005 ($\ge 0.822$; $\ge 0.821$) & 0.841 / 0.824 & 0.866 / 0.870 \\
    \midrule
    \textbf{Kết luận} & \textbf{Đạt cả ba} & \textbf{Đạt cả ba} \\
    \bottomrule
  \end{tabularx}\par
  \vspace{0.6em}
  \begin{itemize}
    \item Tiêu chí lưu vào kho mã nguồn \alert{trước khi huấn luyện}; tập test chỉ chấm \alert{sau khi đã nhận}.
    \item Huấn luyện trên GPU Tesla T4 (Kaggle), tối đa 100 epoch, AdamW; YOLO26s-v6 mất 5.16 giờ.
  \end{itemize}
\end{frame}

\subsection{Giao thức so sánh CNN và Transformer}
\begin{frame}{\stitle}
  \begin{columns}[T]
    \begin{column}{0.48\textwidth}
      \begin{block}{Đồng nhất}
        \begin{itemize}
          \item Cùng dữ liệu v5, cùng file nhãn, cùng khởi tạo COCO.
          \item Ảnh 960, tối đa 100 epoch, dừng sớm sau 20.
          \item Cùng cách chọn checkpoint và cùng đường chấm.
          \item Cùng CPU, cùng frame khi đo tốc độ.
        \end{itemize}
      \end{block}
    \end{column}
    \begin{column}{0.48\textwidth}
      \begin{block}{Giữ theo công thức từng họ}
        Tốc độ học, batch, tăng cường dữ liệu, hàm mất mát. YOLO giữ tỉ lệ ảnh (544$\times$960), RT-DETRv2 kéo giãn
        (960$\times$960), nhiều điểm ảnh hơn 1.8 lần.
      \end{block}
      \begin{block}{Đối chứng}
        Thêm \alert{YOLO26s ở ảnh 1280} (gần cùng số điểm ảnh) để tách tác dụng của kiến trúc và của độ phân giải.
      \end{block}
    \end{column}
  \end{columns}
\end{frame}

\subsection{Pipeline hệ thống}
\begin{frame}{\stitle}
  \centering
  \figorbox{pipeline.pdf}{width=0.92\textwidth,height=0.36\textheight,keepaspectratio}{3cm}
  \vspace{0.3em}
  \begin{columns}[T]
    \begin{column}{0.48\textwidth}
      \begin{itemize}
        \item Phát hiện (ngưỡng 0.1) $\rightarrow$ \alert{lọc box trùng} IoU 0.6 $\rightarrow$ BoT-SORT (không GMC, không
              ReID).
        \item Ánh xạ lớp theo tên, chạy được cả mô hình COCO gốc.
      \end{itemize}
    \end{column}
    \begin{column}{0.48\textwidth}
      \begin{itemize}
        \item \textbf{Bỏ phiếu lớp theo track}: $\hat{c}(k) = \arg\max_c \sum_{t} s_t \mathbf{1}\{c_t = c\}$.
        \item Đếm khi điểm giữa đáy box cắt vạch; xuất video H.264 và CSV.
      \end{itemize}
    \end{column}
  \end{columns}
\end{frame}

% =====================================================================
\section{Kết quả thực nghiệm}
% =====================================================================

\subsection{Phát hiện: test Hà Nội và cảnh mới}
\begin{frame}{\stitle}
  \centering
  \footnotesize
  \begin{tabular}{@{}lcccc@{}}
    \toprule
    & & \multicolumn{3}{c}{\textbf{\texttt{vn\_extra} (camera mới)}} \\
    \cmidrule(lr){3-5}
    \textbf{Mô hình} & \hd{Test Hà Nội}{mAP50} & \hd{Ngày}{mAP50} & \hd{Đêm}{mAP50} & \hd{Đêm}{mAP50-95} \\
    \midrule
    COCO gốc @1280 & \ci{0.686}{0.668--0.705} & \ci{0.458}{0.416--0.502} & \ci{0.437}{0.363--0.498} & 0.308 \\
    v3 & \ci{0.921}{0.911--0.933} & 0.757 & 0.437 & 0.263 \\
    v5 & \ci{\textbf{0.923}}{0.913--0.932} & \ci{0.745}{0.699--0.787} & \ci{0.628}{0.575--0.681} & 0.404 \\
    v6 (bản nhẹ) & \ci{0.921}{0.909--0.932} & \ci{0.760}{0.724--0.795} & \ci{0.755}{0.708--0.795} & 0.495 \\
    \rowcolor{HiRow}
    YOLO26s-v6 & \ci{0.920}{0.911--0.931} & \ci{\textbf{0.830}}{0.801--0.858} & \ci{\textbf{0.887}}{0.861--0.913} &
      \textbf{0.614} \\
    \bottomrule
  \end{tabular}\par
  \vspace{0.4em}
  \begin{itemize}
    \item Tinh chỉnh nâng mAP50 Hà Nội từ 0.686 lên 0.920--0.923; các bản từ v3 trở đi \alert{hoà nhau} ở đây.
    \item Trên camera mới, YOLO26s-v6 cao nhất cả ngày lẫn đêm; frame đêm (v6) nâng đêm từ 0.628 lên 0.755.
  \end{itemize}
\end{frame}

\subsection{Phát hiện: clip Phạm Văn Đồng và tập ngoài}
\begin{frame}{\stitle}
  \begin{columns}[T]
    \begin{column}{0.62\textwidth}
      \centering
      \footnotesize
      \setlength{\tabcolsep}{3.4pt}
      % Tô nền cột mô hình chính như bảng tương ứng trong báo cáo
      \begin{tabular}{@{}lcccc>{\columncolor{HiRow}}c@{}}
        \toprule
        \textbf{Tập} & \hd{COCO}{@1280} & \textbf{v3} & \textbf{v5} & \textbf{v6} & \textbf{YOLO26s-v6} \\
        \midrule
        Clip PVD & 0.664 & 0.891 & 0.892 & \textbf{0.907} & 0.901 \\
        Val UVH-26 & 0.594 & 0.335 & 0.850 & 0.849 & \textbf{0.871} \\
        UA-DETRAC đêm & 0.689 & 0.722 & 0.845 & 0.857 & \textbf{0.867} \\
        UA-DETRAC ngày & 0.690 & 0.709 & 0.812 & 0.852 & \textbf{0.881} \\
        Test Roboflow & 0.255 & -- & 0.834 & 0.844 & \textbf{0.857} \\
        \bottomrule
      \end{tabular}\par
      \vspace{0.5em}
      \begin{itemize}
        \item v3 thấp hơn cả COCO ở UVH-26 (lệch miền Ấn Độ); thêm UVH-26 (v5) đưa lên 0.850.
        \item Frame đêm (v6) không làm kém đi ở camera Ấn Độ, Trung Quốc.
      \end{itemize}
    \end{column}
    \begin{column}{0.34\textwidth}
      \centering
      \figorbox{confusion_matrix.jpg}{width=\linewidth,height=0.48\textheight,keepaspectratio}{3cm}
      \src{Ma trận nhầm lẫn của YOLO26s-v6 trên test Roboflow: nhầm lớp lớn nhất 0.07; lỗi chủ yếu là bỏ sót (18\% xe
        máy).}
    \end{column}
  \end{columns}
\end{frame}

\subsection{So sánh kiến trúc CNN và Transformer}
\begin{frame}{\stitle}
  \centering
  \small
  \begin{tabular}{@{}lcccc@{}}
    \toprule
    \textbf{Tập (mAP50 chính)} & \hd{YOLO26n}{@960} & \hd{YOLO26s}{@960} & \hd{RT-DETRv2-S}{@960$\times$960} &
      \hd{YOLO26s}{@1280} \\
    \midrule
    Test Hà Nội & 0.923 & 0.914 & \textbf{0.929} & 0.915 \\
    \texttt{vn\_extra} ngày & 0.745 & 0.798 & \textbf{0.886} & 0.875 \\
    \texttt{vn\_extra} đêm & 0.628 & 0.713 & \textbf{0.816} & 0.800 \\
    Val UVH-26 & 0.850 & 0.859 & \textbf{0.894} & 0.865 \\
    UA-DETRAC ngày & 0.812 & 0.885 & \textbf{0.924} & 0.885 \\
    \midrule
    mAP50-95 val UVH-26 & 0.696 & 0.721 & \textbf{0.757} & -- \\
    \bottomrule
  \end{tabular}\par
  \vspace{0.5em}
  \begin{enumerate}
    \item Cùng ảnh 960: RT-DETRv2-S bằng hoặc hơn YOLO26s, hơn rõ ở cảnh khó.
    \item Cho YOLO26s \alert{cùng số điểm ảnh} (ảnh 1280): hai họ \alert{hoà nhau} trên dữ liệu Việt Nam, nên phần lớn
          lợi thế đến từ độ phân giải đầu vào.
    \item Trên tập nhãn độc lập, box của RT-DETRv2-S khít hơn khoảng 0.03 mAP50-95.
  \end{enumerate}
\end{frame}

\subsection{Tốc độ trên CPU và GPU}
\begin{frame}{\stitle}
  % Bảng và nhận xét ở trên, đồ thị hai panel ở dưới trải rộng: đặt trong cột hẹp thì chữ trên trục quá nhỏ
  \begin{columns}[T]
    \begin{column}{0.53\textwidth}
      \centering
      \footnotesize
      \setlength{\tabcolsep}{4pt}
      \begin{tabular}{@{}lcc@{}}
        \toprule
        \textbf{Mô hình} & \hd{CPU i5-8250U}{phát hiện (ms)} & \hd{GPU Tesla T4}{pipeline (frame/s)} \\
        \midrule
        YOLO26n & \textbf{100.5} & 42.2 \\
        \rowcolor{HiRow}
        YOLO26s & 261.9 & \textbf{42.6} \\
        YOLO26s @1280 & 455.2 & 30.2 \\
        RT-DETRv2-S & 1526.2 & 16.0\textsuperscript{*} \\
        \bottomrule
      \end{tabular}\par
      \vspace{0.2em}
      \src{* riêng phần phát hiện.}
    \end{column}
    \begin{column}{0.45\textwidth}
      \small
      \begin{itemize}
        \item Trên GPU, YOLO26n và YOLO26s \alert{nhanh như nhau} (khoảng 42 frame/s), nên mô hình chính dùng YOLO26s.
        \item RT-DETRv2-S chậm hơn YOLO26s 5.8 lần trên CPU, khoảng 4 lần trên GPU.
      \end{itemize}
    \end{column}
  \end{columns}
  \vspace{0.1em}
  \centering
  \figorbox{accuracy_vs_speed.pdf}{width=0.86\textwidth,height=0.35\textheight,keepaspectratio}{3cm}
  \src{Độ chính xác theo thời gian phát hiện trên CPU: (a) \texttt{vn\_extra} đêm, (b) val UVH-26.}
\end{frame}

\subsection{Chạy trên laptop với OpenVINO}
\begin{frame}{\stitle}
  \centering
  \small
  % Hai mô hình tách thành hai nhóm như bảng trong báo cáo; in đậm giá trị tốt nhất trong từng nhóm
  \begin{tabular}{@{}lcccc@{}}
    \toprule
    \textbf{Cấu hình, thiết bị} & \hd{Phát hiện}{(ms)} & \hd{Pipeline}{(frame/s)} & \hd{mAP50}{val} & \hd{IDR}{dev} \\
    \midrule
    \multicolumn{5}{@{}l}{\textit{v6 (bản nhẹ, YOLO26n)}} \\
    PyTorch FP32, CPU & 130.0 & 6.6 & 0.914 & \textbf{0.841} \\
    OpenVINO INT8, CPU & 39.9 & 18.2 & \textbf{0.918} & 0.825 \\
    \rowcolor{HiRow}
    OpenVINO FP32, GPU tích hợp & \textbf{34.7} & \textbf{19.5} & 0.914 & 0.825 \\
    \midrule
    \multicolumn{5}{@{}l}{\textit{YOLO26s-v6 (mô hình chính)}} \\
    PyTorch FP32, CPU & 338.0 & 3.2 & 0.911 & 0.866 \\
    OpenVINO FP32, GPU tích hợp & \textbf{80.1} & \textbf{9.8} & \textbf{0.912} & \textbf{0.875} \\
    \bottomrule
  \end{tabular}\par
  \vspace{0.2em}
  \src{In đậm: tốt nhất trong từng mô hình. IDR đo trên VnTraffic, chưa lọc box trùng.}
  \begin{itemize}
    \item Bản nhẹ v6 qua OpenVINO trên GPU Intel UHD 620 đạt \alert{19.5 frame/s}, gấp 3 lần PyTorch.
    \item Ba điều quan trọng: chỉ định thiết bị, tránh tranh chấp luồng CPU, xuất đúng khung 544$\times$960 (nhanh hơn
          ảnh vuông 1.6--2.1 lần).
    \item Đánh đổi: mAP50 hoà PyTorch, IDR thấp hơn 0.016.
  \end{itemize}
\end{frame}

\subsection{Theo dõi: chọn tracker trên dev}
\begin{frame}{\stitle}
  \begin{columns}[T]
    \begin{column}{0.52\textwidth}
      \centering
      \footnotesize
      \begin{tabular}{@{}lcccc@{}}
        \toprule
        \textbf{Tracker (v3)} & \textbf{AssA} & \textbf{IDR} & \textbf{Đổi ID} & \textbf{frame/s} \\
        \midrule
        ByteTrack & 0.615 & 0.708 & 147 & \textbf{4.4} \\
        \rowcolor{HiRow}
        BoT-SORT & \textbf{0.659} & \textbf{0.738} & \textbf{128} & 3.6 \\
        OC-SORT & 0.594 & 0.669 & 206 & 4.0 \\
        \bottomrule
      \end{tabular}\par
      \vspace{0.7em}
      \begin{tabular}{@{}lccc@{}}
        \toprule
        \textbf{Detector} & \textbf{IDR} & \textbf{Đổi ID} & \textbf{Bỏ sót} \\
        \midrule
        v3 & 0.744 & 128 & 2.550 \\
        v5 & 0.827 & \textbf{62} & 1.058 \\
        v6 & 0.841 & 93 & 717 \\
        \rowcolor{HiRow}
        YOLO26s-v6 & \textbf{0.866} & 63 & \textbf{632} \\
        \bottomrule
      \end{tabular}
    \end{column}
    \begin{column}{0.44\textwidth}
      \begin{itemize}
        \item VnTraffic thiếu nhãn xe nên chỉ so bằng \alert{AssA, IDR, số lần đổi ID}.
        \item BoT-SORT tốt nhất; tắt bù chuyển động camera vì camera đứng yên (hoà trên dev, nhanh hơn).
        \item Mô hình phát hiện ảnh hưởng lên theo dõi \alert{nhiều hơn} cấu hình tracker: v3 sang v5 giảm đổi ID
              từ 128 xuống 62.
      \end{itemize}
    \end{column}
  \end{columns}
\end{frame}

\subsection{Giảm số lần đổi ID bằng bộ lọc box trùng}
\begin{frame}{\stitle}
  \centering
  \small
  \begin{tabular}{@{}lcccccc@{}}
    \toprule
    & \multicolumn{2}{c}{\textbf{Dev VnTraffic}} & \multicolumn{2}{c}{\textbf{Dev PVD góc C}} &
      \multicolumn{2}{c}{\textbf{Test PVD góc A}} \\
    \cmidrule(lr){2-3}\cmidrule(lr){4-5}\cmidrule(lr){6-7}
    \textbf{Cấu hình tracker (v6)} & \textbf{Đổi ID} & \textbf{IDR} & \textbf{Bỏ sót} & \textbf{HOTA} & \textbf{Đổi ID} &
      \textbf{HOTA} \\
    \midrule
    Cũ (không lọc) & 93.5 & 0.833 & \textbf{1.521} & 0.689 & 49 & 0.732 \\
    Mở track từ 0.6 & 62.5 & 0.836 & 2.182 & \textbf{0.700} & \textbf{7} & 0.687 \\
    \rowcolor{HiRow}
    Lọc box trùng IoU 0.6 & 63.0 & \textbf{0.841} & 1.535 & 0.695 & 26 & \textbf{0.744} \\
    \bottomrule
  \end{tabular}\par
  \vspace{0.5em}
  \begin{itemize}
    \item YOLO26 không chạy NMS: khoảng 20 trong 82 box mỗi frame là \alert{box phụ} chồng lên cùng một xe.
    \item Nâng ngưỡng mở track đẹp trên VnTraffic nhưng \alert{bỏ sót gần gấp đôi} trên test, vì VnTraffic không có nhãn
          xe nhỏ ở xa. Vì vậy đề tài gán nhãn thêm clip dev góc C và chọn lại trên cả hai tập dev.
    \item Bộ lọc box trùng IoU 0.6: tổng đổi ID trên dev giảm 30\%, bỏ sót chỉ tăng 0.9\%; trên test đổi ID 49 xuống 26.
  \end{itemize}
\end{frame}

\subsection{Theo dõi trên tập test}
\begin{frame}{\stitle}
  \centering
  \small
  \begin{tabular}{@{}lcccccc@{}}
    \toprule
    \textbf{Cấu hình} & \textbf{HOTA} & \textbf{MOTA} & \textbf{IDF1} & \textbf{IDR} & \textbf{IDP} & \textbf{Đổi ID} \\
    \midrule
    \rowcolor{HiRow}
    YOLO26s-v6 (mặc định) & \textbf{0.755} & \textbf{0.874} & \textbf{0.916} & \textbf{0.882} & 0.953 & 22 \\
    YOLO26s-v6, lọc track dưới 3 frame & 0.745 & 0.855 & 0.907 & 0.860 & \textbf{0.959} & \textbf{17} \\
    v6 (bản nhẹ) & 0.744 & 0.847 & 0.905 & 0.859 & 0.955 & 26 \\
    v5, chưa lọc box trùng & 0.718 & 0.813 & 0.880 & 0.822 & 0.947 & 26 \\
    \bottomrule
  \end{tabular}\par
  \vspace{0.6em}
  \begin{itemize}
    \item Clip Phạm Văn Đồng 20 giây, 183 xe, nhãn quỹ đạo đầy đủ.
    % ~\, nới khoảng sau IDF1: nét số 1 của font sát chữ số kế tiếp, khoảng thường đọc thành "IDF10.916"
    \item Mô hình chính: \alert{HOTA 0.755, MOTA 0.874, IDF1~\,0.916}, đổi ID 22 lần.
    \item IDP cao hơn IDR: ít gán nhầm danh tính, phần mất chủ yếu do bỏ sót (1.353 box sót, 340 box thừa).
    \item Với cùng v5, BoT-SORT cũng tốt nhất trên test (ByteTrack HOTA 0.698, OC-SORT 0.692).
  \end{itemize}
\end{frame}

\subsection{Ứng dụng: sai số đếm xe qua vạch}
\begin{frame}{\stitle}
  \centering
  \small
  \begin{tabular}{@{}lccccc@{}}
    \toprule
    & \multicolumn{2}{c}{\textbf{Ban ngày (73 lượt)}} & \multicolumn{3}{c}{\textbf{Ban đêm (176 lượt, đếm tay)}} \\
    \cmidrule(lr){2-3}\cmidrule(lr){4-6}
    \textbf{Mô hình} & \textbf{Đếm (vào / ra)} & \textbf{Sai số} & \textbf{Sót} & \textbf{Thừa} & \textbf{Sai số} \\
    \midrule
    v5 & 0 / 72 & 3 (4.1\%) & 15 & \textbf{0} & 15 (8.5\%) \\
    YOLO26s (dữ liệu v5) & 0 / 71 & \textbf{2 (2.7\%)} & 7 & 5 & 8 (4.5\%) \\
    v6 (bản nhẹ) & 0 / 71 & \textbf{2 (2.7\%)} & \textbf{3} & 2 & \textbf{3 (1.7\%)} \\
    \rowcolor{HiRow}
    YOLO26s-v6 & 0 / 70 & 3 (4.1\%) & \textbf{3} & 2 & \textbf{3 (1.7\%)} \\
    \bottomrule
  \end{tabular}\par
  \vspace{0.5em}
  \begin{itemize}
    \item Ban ngày: phần lớn sai số là hiện tượng của nhãn ở tán cây; lỗi thật là một Ford Transit đếm thành ô tô.
    \item Ban đêm: mọi lỗi ở xe máy, sót trong \alert{cụm xe đi sát nhau}; thừa do box dao động quanh vạch.
    \item Học thêm cảnh đêm đưa sai số đêm từ 8.5\% xuống \alert{1.7\%}.
  \end{itemize}
\end{frame}

\subsection{Phân tích lỗi: bỏ sót ban ngày}
\begin{frame}{\stitle}
  \begin{columns}[T]
    \begin{column}{0.49\textwidth}
      \centering
      \figorbox{error_1.jpg}{width=\linewidth,height=0.42\textheight,keepaspectratio}{3cm}
      \src{Xe máy nhỏ ở xa bị sót (\texttt{vn\_extra} ngày).}
    \end{column}
    \begin{column}{0.49\textwidth}
      \centering
      \figorbox{error_2.jpg}{width=\linewidth,height=0.42\textheight,keepaspectratio}{3cm}
      \src{Cụm xe máy chờ đèn đỏ bị sót (test Hà Nội).}
    \end{column}
  \end{columns}
  \vspace{0.4em}
  \begin{itemize}
    \item Đỏ là bỏ sót, xanh lá là đúng, tím là thừa, xám là vùng bỏ qua.
    \item Hai dạng bỏ sót chính: xe quá nhỏ sau khi thu ảnh về 960 điểm ảnh, và xe máy dừng sát nhau che lẫn nhau.
  \end{itemize}
\end{frame}

\subsection{Phân tích lỗi: ban đêm và đổi ID}
\begin{frame}{\stitle}
  \begin{columns}[T]
    \begin{column}{0.44\textwidth}
      \centering
      \figorbox{error_3.jpg}{width=\linewidth,height=0.40\textheight,keepaspectratio}{3cm}
      \src{Ô tô ban đêm chỉ thấy đèn bị sót.}
    \end{column}
    \begin{column}{0.54\textwidth}
      \centering
      \figorbox{error_7.jpg}{width=\linewidth,height=0.30\textheight,keepaspectratio}{2.4cm}
      \src{Đổi ID khi xe bị cây che: ID 296 (frame 340) sang ID 780 (frame 341).}
    \end{column}
  \end{columns}
  \vspace{0.4em}
  \begin{itemize}
    \item Ban đêm: xe máy chở hàng có đèn chói bị gọi là ô tô; mô hình hay chỉ khoanh phần đèn hậu nên một xe vừa bị
          tính sót vừa sinh box thừa.
    \item Đổi ID: xe đi sau cây, box co lại theo phần nhìn thấy, tracker tạo track mới. Hướng sửa: \alert{ReID riêng cho
          xe}.
  \end{itemize}
\end{frame}

\subsection{Độ tin cậy của kết luận}
\begin{frame}{\stitle}
  \begin{columns}[T]
    \begin{column}{0.48\textwidth}
      \begin{block}{Huấn luyện lặp lại với seed khác}
        \begin{itemize}
          \item v5 và YOLO26s ba lần, v6 và YOLO26s-v6 hai lần.
          \item YOLO26s hơn YOLO26n ở \texttt{vn\_extra} ngày và đêm \alert{ở mọi seed}.
          \item Đêm dao động lớn: YOLO26n 0.510--0.638, YOLO26s 0.713--0.821.
          \item Lần kém hơn của YOLO26s-v6 (0.847) vẫn hơn lần tốt nhất của YOLO26s học dữ liệu v5 (0.821).
        \end{itemize}
      \end{block}
    \end{column}
    \begin{column}{0.48\textwidth}
      \begin{block}{Thiên lệch và dùng lại tập test}
        \begin{itemize}
          \item Box sơ bộ giữ nguyên mang dấu mô hình tạo ra nó, nên ở test Hà Nội và clip PVD chỉ dùng mAP50.
          \item Khoảng tin cậy lấy mẫu theo khối 10 giây: kết luận chính giữ nguyên.
          \item Một lần kết quả test làm loại cấu hình chọn trên dev; đã ghi rõ trong báo cáo.
        \end{itemize}
      \end{block}
    \end{column}
  \end{columns}
\end{frame}

\subsection{Web demo}
\begin{frame}{\stitle}
  % Hai ảnh chụp màn hình nằm ngang (2 : 1) đặt cạnh nhau ở trên cho to, mô tả ở dưới
  \begin{columns}[T]
    \begin{column}{0.49\textwidth}
      \centering
      \fbox{\includegraphics[width=0.96\linewidth]{app_gpu_demo.jpg}}
      \src{Video kết quả của clip Phạm Văn Đồng.}
    \end{column}
    \begin{column}{0.49\textwidth}
      \centering
      \fbox{\includegraphics[width=0.96\linewidth]{app_counts.jpg}}
      \src{Bảng đếm theo lớp và hướng, biểu đồ cộng dồn.}
    \end{column}
  \end{columns}
  \vspace{0.5em}
  \begin{itemize}
    \item Streamlit: chọn clip hoặc tải video, chọn mô hình, đặt vạch đếm.
    \item Kết quả: video có box, ID, vết di chuyển; bảng đếm theo lớp và hướng; biểu đồ cộng dồn; tải CSV.
    \item Chế độ trình chiếu nạp video xử lý sẵn bằng \alert{YOLO26s-v6}; chế độ chạy trên laptop dùng bản nhẹ v6 OpenVINO.
  \end{itemize}
\end{frame}

% =====================================================================
\section{Kết luận và hướng phát triển}
% =====================================================================

\subsection{Kết quả chính}
\begin{frame}{\stitle}
  \begin{itemize}
    \item \textbf{Phát hiện}: YOLO26s-v6 trên camera mới đạt mAP50 0.830 ban ngày, \alert{0.887 ban đêm} (COCO gốc 0.458
          và 0.437).
    \item \textbf{Theo dõi}: trên clip test HOTA 0.755, MOTA 0.874, IDF1~\,0.916; đổi ID 22 lần trên 183 xe.
    \item \textbf{Đếm xe}: sai số 4.1\% ban ngày (73 lượt), \alert{1.7\%} ban đêm (176 lượt đếm tay).
    \item \textbf{Kiến trúc}: RT-DETRv2-S chính xác hơn YOLO26s ở cùng ảnh 960 nhưng hoà khi YOLO nhìn cùng số điểm ảnh,
          và chậm hơn 4--5.8 lần.
    \item \textbf{Triển khai}: khoảng 43 frame/s trên GPU T4; laptop không GPU rời dùng bản nhẹ v6 OpenVINO 19.5 frame/s.
    \item \textbf{Dữ liệu đúng cảnh} có tác dụng lớn hơn tăng kích thước ảnh hay đổi kiến trúc.
  \end{itemize}
\end{frame}

\subsection{Đóng góp và hạn chế}
\begin{frame}{\stitle}
  \begin{columns}[T]
    \begin{column}{0.48\textwidth}
      \begin{block}{Đóng góp}
        \begin{itemize}
          \item Tập test Việt Nam có người kiểm, gồm cảnh đêm và nhãn quỹ đạo.
          \item Chọn trên dev theo tiêu chí viết trước; tách camera đã học và camera mới.
          \item So sánh CNN và Transformer có đối chứng cùng số điểm ảnh.
          \item Chạy được trên cả máy chủ GPU lẫn laptop.
        \end{itemize}
      \end{block}
    \end{column}
    \begin{column}{0.48\textwidth}
      \begin{block}{Hạn chế}
        \begin{itemize}
          \item Tập test nhỏ (20--30 frame), một clip theo dõi 20 giây, ba địa điểm.
          \item Chưa có nhãn theo dõi cho cảnh đêm.
          \item Nhãn train đêm là nhãn tự động, thiếu xe nhỏ ở xa.
          \item Chưa có lớp xe van; tốc độ T4 mới đo một luồng FP32.
        \end{itemize}
      \end{block}
    \end{column}
  \end{columns}
\end{frame}

\subsection{Hướng phát triển}
\begin{frame}{\stitle}
  \begin{itemize}
    \item Thêm dữ liệu đúng góc camera giao thông Việt Nam, chia theo camera; gán nhãn theo dõi cho cảnh đêm.
    \item Huấn luyện RT-DETRv2-S trên dữ liệu v6, lặp lại vài lần để so ban đêm cùng dữ liệu.
    \item Chỉ đếm khi xe ở hẳn bên kia vạch vài frame, chặn đếm lặp do box dao động.
    \item Mô hình ReID riêng cho xe để giảm đổi ID khi bị che.
    \item TensorRT, FP16 trên GPU; đo số camera một GPU xử lý kịp cùng lúc.
    \item Tách khâu vẽ và ghi video sang luồng riêng hoặc dùng bộ mã hoá phần cứng.
  \end{itemize}
\end{frame}

% =====================================================================
\bibslide
\qaslide{Em xin chân thành cảm ơn Quý Thầy Cô đã lắng nghe.}

\end{document}
